\paragraph{Control for source covariance.}
To test whether the source benefit is explained by its covariance, we compare
against a single Gaussian with approximately the same covariance as the harmonic
mixture. We estimate that covariance from 128 trees and average over permutations
of identical atom species. An independent 8,192-tree estimate gives a relative
Frobenius error of 1.85--6.25\% across the 22 compositions. We match the training
data, optimizer, network, and sampling budgets. This follow-up compares the new
Gaussian models with the already frozen harmonic checkpoints.

\begin{table}[t]
\centering
\caption{Covariance-matched Gaussian versus harmonic source. Entries count
valid outputs over all attempts. Training and inference budgets are matched;
the Gaussian uses the calibrated covariance estimate from 128 trees.}
\begin{tabular}{llrr}
\toprule
Set & Run & Covariance Gaussian & Harmonic tree\\
\midrule
Development & 1 &332/768&436/768\\
            & 2 &313/768&388/768\\
Additional  & 1 &286/640&413/640\\
            & 2 &305/640&375/640\\
\bottomrule
\end{tabular}\label{tab:covariance}
\end{table}

Harmonic improves graph validity over this Gaussian by 11.65 percentage points
on development (paired 95\% interval [8.59,14.78]; descriptive composition interval
[8.07,15.43]) and 15.39 points on the additional set (paired [11.64,19.06];
composition [11.95,18.98]). The difference is positive in both runs on both sets.
The covariance-matched Gaussian also performs worse than the isotropic Gaussian.
These results support a benefit beyond this calibrated second-moment approximation,
consistent with a contribution from the mixture structure. The comparison adds
2,816 outputs while holding parameter count, initial self-conditioning weights,
training structures, and optimization fixed.
